% Pursuing Stacks (cotes 134-2 to 134-8): archive cote 134-7, site page 28 (PDF page 29), typescript page 470.
% typescript. Draft compared: ps6.tex:1298-1359.
% AI-transcribed by Claude Opus 5.5 (claude-opus-5-5, adaptive thinking, effort high), Message Batches API, batch ps3 (msgbatch_01XdFmuVVBAcgc7EFxdtCqsj), 2026-11-12; not yet checked by eye.
\scanpage{134-7/28}
\mspage{470}

\deleted{from} the Kan condition that $\pi_1$ may be interpreted as a quotient set of \xout{X}
\inserted{$E \overset{\mathrm{def}}{=}$}$X_1(k)$ (without having to pass to the full free group generated by this
set), with a set of relations
\[ \text{(F)} \qquad z_i = x_iy_i \quad , \quad i \text{ in } X_2(k)=I, \]
indexed by $X_2(k)$, where
\[ i\mapsto x_i \ , \ i\mapsto y_i \ , \ i\mapsto z_i \]
are the three boundary maps, remembering moreover \inserted{the Kan condition} that the three maps
\[ \diagram{\begin{tikzcd}[cramped,sep=small] I \ar[r]\ar[r,shift left=1.5]\ar[r,shift right=1.5] & E\mathrm{x}E \end{tikzcd}} \ , \quad i\mapsto(x_i,y_i) \ , \quad i\mapsto(x_i,z_i) \ , \quad i\mapsto(y_i,z_i) \]
are \uline{surjective} , which implies indeed that \xout{xxx} any element of the group \inserted{$\pi$}
described by the set of generators E and relations (F) \xout{xx} comes from an
element in E. Replacing k by any k-algebra k', we see that we have a
presheaf\xout{\illegible{1}}
\[ k'\mapsto \pi_1(X_*(k')) = \pi_1(X_*\otimes_k k') \]
on the category $\mathrm{Aff}_{/k}$ of affine schemes over k, with values in the
category of groups, which may be viewed (as a presheaf of sets) as a
quotient presheaf of the presheaf on $\mathrm{Aff}_{/k}$ defined by $X_1$ . \xout{\unsure{In order}} We
feel that this \xout{\illegible{1}}\inserted[typed]{presheaf}will fit into a reasonable ``schematic'' set up,
only if it turns out to be a \uline{sheaf}, and \xout{even} more exactingly still, if
this sheaf is \xout{representable}\inserted[typed]{isomorphic} (as a sheaf of sets) to one stemming from an
object of U(k) i.e. if it is isomorphic to a sheaf W(M), for suitable
k-module M (not necessarily a projective one). If we denote by G this
object of U(k), it will be endowed with a group structure, and \inserted{it is} \uline{this}
\uline{group object of} U(k), \uline{rather than just the} \inserted{set-theoretic} \uline{group of it's sections} i.e.
\uline{of k-valued ``points''}, \inserted{which} \uline{merits to be viewed as the ``true''} $\pi_1(X_*)$. To say
it differently, whereas the higher $\pi_i(X_*) \overset{\mathrm{def}}{=} \pi_i(X_*(k))$ \xout{xxx} (for $i\geq 2$),
in the cases considered so far, should be viewed as being not mere abeli\xout{a}
\hand{-}an groups, but moreover endowed with a natural \over{k}{k}-module structure,
in the case when i=1 i.e. for the fundamental group $\pi_1(X_*)$, the natu-%
ral structure to expect on this (possibly non commutative) group is a
\inserted{``unipotent} ``schematic'' structure, namely essentially a \inserted{pointed} ``parametrisation'' of this
group by elements of a suitable k-module $M_1$ , in such a way that the
composition law is expressed in terms of a polynomial law, making sense
therefore not only for k-valued points i.e. \inserted{for} elements of M, but \inserted{for} k'-valued
points as well \hand{(}for any k-algebra k'\hand{),} namely defining a group law on

\margin{28}
